\documentclass[10pt,a4paper]{amsart}
\usepackage[margin=19mm]{geometry}
\usepackage{amsmath,amssymb,mathtools}
\usepackage[scr=rsfs,cal=euler]{mathalfa}
\usepackage{xfrac}
\usepackage[unicode,hidelinks,bookmarks=false]{hyperref}
\newcommand{\ie}{i.e.\ }
\newcommand{\cf}{cf.\ }
\newcommand{\resp}{respectively\ }
\newcommand{\Subsection}{\subsection}
\providecommand{\nobreakdash}{\nobreak\hbox{-}}
\setlength{\emergencystretch}{2em}
\multlinegap0pt
\usepackage{bm}
\usepackage{bbm}
\usepackage{dsfont}
\usepackage{xspace}
\usepackage{cancel}
\usepackage{mathtools}
\usepackage{amsthm}
\makeatletter
\@ifundefined{lemma}{\newtheorem{lemma}{Lemma}}{}
\makeatother
\newcommand{\as}{\overset{\mathrm{a.s.}}{\longrightarrow}}
\newcommand{\pconv}{\overset{{p}}{\to}}
\newcommand{\prox}{\mathrm{prox}}
\newcommand{\pr}{\mathbb{P}}
\renewcommand{\hat}{\widehat}
\title[Moment-Based Adjustments of Statistical Inference in High-Di]{Moment-Based Adjustments of Statistical Inference in High-Dimensional Generalized Linear Models}
\author{Kazuma Sawaya, Yoshimasa Uematsu, Masaaki Imaizumi}
\date{Source: arXiv:2305.17731v5 (CC BY 4.0)}
\begin{document}
\maketitle

\noindent\emph{Source and licence.} Moment-Based Adjustments of Statistical Inference in High-Dimensional Generalized Linear Models. Version arXiv:2305.17731v5 (CC BY 4.0), distributed under the Creative Commons Attribution 4.0 International licence, as recorded in the arXiv licence metadata for this exact version and shown on the abstract page https://arxiv.org/abs/2305.17731v5. Full rights notice: licenses/arxiv-2305.17731-CC-BY-4.0.txt.\par\smallskip

\noindent\textit{Editorial scope.} Counted unit: the Lemma whose source label is \texttt{lem:SE-conv} in the article (source0.tex lines 2717--2762, source label \texttt{lem:SE-conv}), delivered together with the article's own complete original proof of it (source0.tex lines 2764--3163, one \texttt{proof} environment of 400 lines). No mathematics is written, repaired or re-derived: statement and proof are the article's own text line for line. Required context retained uncounted: none. Every definition, display and label of those lines is retained unchanged. Omitted from this package and not redistributed: 1--2716, 2763--2763, 3164--3300. Nothing inside a retained excerpt is omitted or altered: every retained body is delivered exactly as the article's lines are written, so no omission receipt of any kind accompanies this package. Citations: the retained excerpts carry no citation command at all, so no bibliography entry of the article is transcribed and no reference is redistributed. Reference adaptation: none was needed and none was made; every reference inside the delivered unit resolves inside it, so no unresolved reference or citation is produced. Printed slips kept literal: no printed word, symbol, hypothesis or number of the counted statement or of its proof was altered. Layout and class: the article's own class is \texttt{amsart} with packages including amsthm, amsmath, amssymb, amsfonts, amsaddr, natbib, txfonts, fullpage, inputenc, fontenc, hyperref, url, booktabs, nicefrac; the wrapper is the lane's pinned \texttt{amsart} editorial wrapper at 10pt on A4, which restates the theorem-like environments the retained text needs and adds the article's own macro definitions that the retained text needs (\texttt{\textbackslash as}, \texttt{\textbackslash hat}, \texttt{\textbackslash pconv}, \texttt{\textbackslash pr}, \texttt{\textbackslash prox}), with extra wrapper packages (bm, bbm, dsfont, xspace, cancel, mathtools). The delivered numbering is the unit's own. Assets: no retained span contains a figure environment, a graphics-inclusion command, a PDF-page inclusion, a table or a TikZ picture; no image file is copied into this package, the package declares no asset dependency and no image file is redistributed with it. Licence: version arXiv:2305.17731v5 of this article is distributed under the Creative Commons Attribution 4.0 International licence, as recorded in the arXiv licence metadata for this exact version and shown on the abstract page https://arxiv.org/abs/2305.17731v5. A full rights notice is delivered at licenses/arxiv-2305.17731-CC-BY-4.0.txt. Characters: every retained excerpt is pure ASCII, so the delivered engine, XeLaTeX with its default Latin Modern text font, reports no missing character.
\begin{lemma}
\label{lem:SE-conv}
Let $Q_1,Q_2\sim\mathcal{N}(0,1)$ be independent, and let
$\bar\varepsilon$ be independent of $(Q_1,Q_2)$. Suppose that
$g:\mathbb{R}\to\mathbb{R}$ is continuously differentiable and
nondecreasing, and let $G:\mathbb{R}\to\mathbb{R}$ satisfy $G'=g$.
Suppose further that
$h:\mathbb{R}^2\to\mathbb{R}$ is Lipschitz continuous in its first
argument and that $\mathbb{E}\left[h(0,\bar\varepsilon)^2\right]<\infty.$

Let $\gamma>0$, and let
$\mathcal I_\gamma\subset(0,\infty)$ be a compact interval such that $\gamma\in\operatorname{int}(\mathcal I_\gamma).$
Let $\Theta\subset\mathbb{R}\times(0,\infty)^2$ be compact. 
Define the state-evolution residual map $\mathcal F:
\Theta\times\mathcal I_\gamma
\rightarrow
\mathbb{R}^3$.
Let $\bm\vartheta_0:=(\mu,\sigma^2,\eta)\in\Theta$ be the unique solution in $\Theta$ to the oracle state-evolution system
\[
\mathcal F(\bm\vartheta,\gamma)=\bm 0.
\]
Let $\hat\gamma\ge0$ satisfies
\[
\hat\gamma^2\pconv\gamma^2,
\]
and let $\hat{\bm\vartheta}:=(\hat\mu,\hat\sigma^2,\hat\eta)$ be a measurable random vector satisfying
\begin{align}
\pr\left(
\hat\gamma\in\mathcal I_\gamma,\,
\hat{\bm\vartheta}\in\Theta
\right)
&\rightarrow1,
\label{eq:plugin-SE-compactness}\\
\left\|
\mathcal F(\hat{\bm\vartheta},\hat\gamma)
\right\|
&\pconv0.
\label{eq:plugin-SE-approx-root}
\end{align}
Then, we obtain
\[
\hat{\bm\vartheta}
\pconv
\bm\vartheta_0.
\]
\end{lemma}

\begin{proof}
For $\bm\vartheta=(m,v,a)\in\Theta$ and $\varsigma\in\mathcal I_\gamma$, we can write
\begin{align}
\mathcal F(\bm\vartheta,\varsigma)
:=
\begin{pmatrix}
\kappa^2v
-
a^2\mathbb{E}\left[
R_{\bm\vartheta,\varsigma}^2
\right]
\\[1mm]
\mathbb{E}\left[
\varsigma Q_1R_{\bm\vartheta,\varsigma}
\right]
\\[1mm]
1-\kappa
-
\mathbb{E}\left[
\left\{
1+a g'(D_{\bm\vartheta,\varsigma})
\right\}^{-1}
\right]
\end{pmatrix},
\label{eq:SE-residual-map}
\end{align}
where
\begin{align}
\bar Y_{\varsigma}
&:=
h(\varsigma Q_1,\bar\varepsilon),
\qquad S_{\bm\vartheta,\varsigma}
:=
m\varsigma Q_1+\sqrt{\kappa v}\,Q_2,
\label{eq:SE-linear-varsigma}\\
D_{\bm\vartheta,\varsigma}
&:=
{\prox}_{aG}
\left(
S_{\bm\vartheta,\varsigma}
+a\bar Y_{\varsigma}
\right),
\qquad R_{\bm\vartheta,\varsigma}
:=
\bar Y_{\varsigma}
-
g(D_{\bm\vartheta,\varsigma}).
\label{eq:SE-residual-varsigma}
\end{align}

\textbf{Step 1: Continuity}.
We first show that $(\bm\vartheta,\varsigma)
\mapsto
\mathcal F(\bm\vartheta,\varsigma)$
is continuous on
$\Theta\times\mathcal I_\gamma$.

Consider an arbitrary convergent sequence $(\bm\vartheta_\ell,\varsigma_\ell)
\rightarrow
(\bm\vartheta,\varsigma)$ in $\Theta\times\mathcal I_\gamma,$
where $\bm\vartheta_\ell=(m_\ell,v_\ell,a_\ell)$ and $\bm\vartheta=(m,v,a).$
For brevity, write
\[
\bar Y_\ell
:=
\bar Y_{\varsigma_\ell},
\qquad
S_\ell
:=
S_{\bm\vartheta_\ell,\varsigma_\ell},
\qquad
D_\ell
:=
D_{\bm\vartheta_\ell,\varsigma_\ell},
\qquad
R_\ell
:=
R_{\bm\vartheta_\ell,\varsigma_\ell},
\]
and define
$\bar Y$, $S$, $D$, and $R$ analogously at
$(\bm\vartheta,\varsigma)$.

Since $h$ is Lipschitz continuous in its first argument,
\[
\bar Y_\ell
=h(\varsigma_\ell Q_1,\bar\varepsilon)
\as
h(\varsigma Q_1,\bar\varepsilon)
=\bar Y.
\]
Moreover,
\[
S_\ell
=m_\ell\varsigma_\ell Q_1+\sqrt{\kappa v_\ell}\,Q_2
\as
m\varsigma Q_1+\sqrt{\kappa v}\,Q_2=
S.
\]

By the first-order condition for the proximal operator,
\begin{align}
D_\ell+a_\ell g(D_\ell)
&=
S_\ell+a_\ell\bar Y_\ell,
\label{eq:prox-FOC-seq}\\
D+a g(D)
&=
S+a\bar Y.
\label{eq:prox-FOC-limit}
\end{align}
Subtracting \eqref{eq:prox-FOC-limit} from
\eqref{eq:prox-FOC-seq} gives
\begin{align}
&D_\ell-D
+
a_\ell\{g(D_\ell)-g(D)\}
\notag\\
&\qquad=
S_\ell+a_\ell\bar Y_\ell
-
S-a\bar Y
+
(a-a_\ell)g(D).
\label{eq:prox-continuity-decomp}
\end{align}
Since $g$ is nondecreasing,
\[
(D_\ell-D)\{g(D_\ell)-g(D)\}\ge0.
\]
Consequently, the two terms on the left-hand side of
\eqref{eq:prox-continuity-decomp} have the same sign, and hence
\begin{align}
|D_\ell-D|
&\le
\left|
S_\ell+a_\ell\bar Y_\ell-S-a\bar Y\right|
+|a_\ell-a|\,|g(D)|
\as0.
\label{eq:prox-joint-continuity}
\end{align}
Also, it follows from the continuity of $g$ and $g'$ that
\begin{align}
R_\ell
&\as R,
\label{eq:residual-pointwise-conv}\\
\left\{
1+a_\ell g'(D_\ell)
\right\}^{-1}
&\as
\left\{
1+a g'(D)
\right\}^{-1}.
\label{eq:derivative-pointwise-conv}
\end{align}

We next construct integrable envelopes in order to pass the limits through
the expectations in \eqref{eq:SE-residual-map}. Since
$\Theta\times\mathcal I_\gamma$ is compact, there exist finite constants
$M,\bar v,\bar a,\bar\gamma>0$ and $\underline a>0$ such that
\[
|m|\le M,\qquad
v\le\bar v,\qquad
\underline a\le a\le\bar a,\qquad
\varsigma\le\bar\gamma
\]
for every
$(m,v,a)\in\Theta$ and
$\varsigma\in\mathcal I_\gamma$. Therefore, for a constant
$C_1>0$ independent of
$(\bm\vartheta,\varsigma)$,
\begin{align}
|S_{\bm\vartheta,\varsigma}|
\le
C_1(|Q_1|+|Q_2|).
\label{eq:linear-envelope}
\end{align}

For $a>0$, define $H_a(t):=t+a g(t).$
Since $g$ is nondecreasing, $H_a$ is strictly increasing and $|H_a(t)-H_a(s)|\ge |t-s|$ for all $s,t\in\mathbb{R}.$
Thus, $H_a^{-1}$ is $1$-Lipschitz. Since
\[
D_{\bm\vartheta,\varsigma}
=
H_a^{-1}\left(S_{\bm\vartheta,\varsigma}+a\bar Y_{\varsigma}\right)
\]
and
\[
H_a^{-1}(a g(0))=0,
\]
we obtain
\begin{align}
|D_{\bm\vartheta,\varsigma}|
&=|H_a^{-1}\left(S_{\bm\vartheta,\varsigma}+a\bar Y_{\varsigma}\right)-H_a^{-1}\left(ag(0)\right)|+|H_a^{-1}\left(ag(0)\right)| \\
&\le
\left|
S_{\bm\vartheta,\varsigma}+a\bar Y_{\varsigma}-a g(0)
\right|.
\label{eq:prox-envelope}
\end{align}
Furthermore, the proximal first-order condition implies
\[
D_{\bm\vartheta,\varsigma}-S_{\bm\vartheta,\varsigma}
=aR_{\bm\vartheta,\varsigma}.
\]
Combining this identity with \eqref{eq:linear-envelope} and
\eqref{eq:prox-envelope} yields, for a constant $C_2>0$,
\begin{align}
|R_{\bm\vartheta,\varsigma}|
\le
C_2
\left\{
1+|Q_1|+|Q_2|+|\bar Y_{\varsigma}|
\right\}
\label{eq:residual-envelope}
\end{align}
uniformly over
$(\bm\vartheta,\varsigma)
\in\Theta\times\mathcal I_\gamma$.

Let $L_h$ be a Lipschitz constant of $h$ with respect to its first
argument. Then
\begin{align}
\sup_{\varsigma\in\mathcal I_\gamma}
|\bar Y_{\varsigma}|
&\le
|h(0,\bar\varepsilon)|
+
L_h\bar\gamma |Q_1|.
\label{eq:Y-envelope}
\end{align}
By assumption, $\mathbb{E}\left[
h(0,\bar\varepsilon)^2
\right]<\infty.$
Hence, \eqref{eq:residual-envelope} and \eqref{eq:Y-envelope} imply that
\[
\sup_{\substack{
\bm\vartheta\in\Theta\\
\varsigma\in\mathcal I_\gamma
}}
|R_{\bm\vartheta,\varsigma}|^2
\]
is bounded by an integrable random variable. Similarly,
\[
\sup_{\substack{
\bm\vartheta\in\Theta\\
\varsigma\in\mathcal I_\gamma
}}
\left|
\varsigma Q_1R_{\bm\vartheta,\varsigma}
\right|
\]
is bounded by an integrable random variable, by the Cauchy--Schwarz
inequality. Finally, since $g'\ge0$,
\[
0<
\left\{
1+a g'(D_{\bm\vartheta,\varsigma})
\right\}^{-1}
\le1.
\]
Therefore, \eqref{eq:residual-pointwise-conv},
\eqref{eq:derivative-pointwise-conv}, and the dominated convergence theorem
show that every component of
$\mathcal F(\bm\vartheta,\varsigma)$ is continuous on
$\Theta\times\mathcal I_\gamma$.

\textbf{Step 2: Consistency}. 
We now prove the consistency of the plug-in solution. Since $\hat\gamma^2\pconv\gamma^2$, $\hat\gamma\ge0$, and $\gamma>0,$
the continuous mapping theorem gives
$\hat\gamma\pconv\gamma.$

Fix any $\epsilon>0$, and define
\[
\Theta_\epsilon
:=
\left\{
\bm\vartheta\in\Theta:
\|\bm\vartheta-\bm\vartheta_0\|\ge\epsilon
\right\}.
\]
If $\Theta_\epsilon=\varnothing$, the conclusion is immediate. Otherwise,
$\Theta_\epsilon$ is compact.

By assumption, $\bm\vartheta_0$ is the unique zero in $\Theta$ of $\bm\vartheta
\mapsto
\mathcal F(\bm\vartheta,\gamma).$
Hence, $\|\mathcal F(\bm\vartheta,\gamma)\|>0$ for every $\bm\vartheta\in\Theta_\epsilon.$
Since
$\bm\vartheta\mapsto
\|\mathcal F(\bm\vartheta,\gamma)\|$
is continuous and $\Theta_\epsilon$ is compact, it attains a strictly
positive minimum:
\begin{align}
c_\epsilon
:=
\inf_{\bm\vartheta\in\Theta_\epsilon}
\|\mathcal F(\bm\vartheta,\gamma)\|
>0.
\label{eq:SE-root-separation}
\end{align}

Because $\mathcal F$ is continuous on the compact set
$\Theta\times\mathcal I_\gamma$, it is uniformly continuous there.
Therefore, there exists $\delta_\epsilon^{\rm uc}>0$ such that
\[
|\varsigma-\gamma|<\delta_\epsilon
\quad\implies\quad
\sup_{\bm\vartheta\in\Theta}
\left\|
\mathcal F(\bm\vartheta,\varsigma)
-
\mathcal F(\bm\vartheta,\gamma)
\right\|
<
\frac{c_\epsilon}{3}.
\]
We choose positive $\delta_\epsilon\le\delta_\epsilon^{\rm uc}$ sufficiently small that $(\gamma-\delta_\epsilon,\gamma+\delta_\epsilon)
\subset\mathcal I_\gamma.$

Define the event
\begin{align}
A_{n,\epsilon}
:=
&
\left\{
\hat\gamma\in\mathcal I_\gamma
\right\}
\cap
\left\{
\hat{\bm\vartheta}\in\Theta
\right\}
\cap
\left\{
|\hat\gamma-\gamma|<\delta_\epsilon
\right\}
\cap
\left\{
\left\|
\mathcal F(\hat{\bm\vartheta},\hat\gamma)
\right\|
<
\frac{c_\epsilon}{3}
\right\}.
\label{eq:SE-consistency-event}
\end{align}
By \eqref{eq:plugin-SE-compactness},
\eqref{eq:plugin-SE-approx-root}, and
$\hat\gamma\pconv\gamma$, we have $\pr(A_{n,\epsilon})\rightarrow1.$

Suppose that $A_{n,\epsilon}$ occurs and that $\|\hat{\bm\vartheta}-\bm\vartheta_0\|
\ge\epsilon.$
Then we have
$\hat{\bm\vartheta}\in\Theta_\epsilon$, so that
\eqref{eq:SE-root-separation} gives
\[
c_\epsilon
\le
\left\|
\mathcal F(\hat{\bm\vartheta},\gamma)
\right\|.
\]
On the other hand, by the triangle inequality,
\begin{align*}
\left\|
\mathcal F(\hat{\bm\vartheta},\gamma)
\right\|
&\le
\left\|
\mathcal F(\hat{\bm\vartheta},\gamma)
-
\mathcal F(\hat{\bm\vartheta},\hat\gamma)
\right\|
+
\left\|
\mathcal F(\hat{\bm\vartheta},\hat\gamma)
\right\|\\
&<
\frac{c_\epsilon}{3}
+
\frac{c_\epsilon}{3}
<
c_\epsilon,
\end{align*}
which is a contradiction. Therefore, on $A_{n,\epsilon}$,
\[
\|\hat{\bm\vartheta}-\bm\vartheta_0\|<\epsilon.
\]
Consequently,
\[
\pr\left(
\|\hat{\bm\vartheta}-\bm\vartheta_0\|
\ge\epsilon
\right)
\le
\pr(A_{n,\epsilon}^c)
\rightarrow0.
\]
This completes the proof.
\end{proof}

\end{document}
